% !TEX program = xelatex
% ============================================================
%  Arabic Professional Resume Template — Nibras Center
%  Compile with: xelatex main.tex (run twice)
%  Engine: XeLaTeX (required for polyglossia + bidi + fontspec)
% ============================================================

\documentclass[10pt,a4paper]{article}

% ---- Geometry ----
\usepackage[a4paper,margin=1.5cm,top=1.5cm,bottom=1.5cm]{geometry}

% ---- Core packages (order matters for bidi) ----
\usepackage{bidi}
\usepackage[no-math]{fontspec}
\usepackage[quiet]{polyglossia}

% ---- Fonts ----
\setmainfont[Scale=0.85]{NotoKufiArabic-Regular.ttf}
\newfontfamily\arabicfont[Script=Arabic,Scale=0.85]{NotoKufiArabic-Regular.ttf}
\newfontfamily\englishfont[Scale=0.8]{TeX Gyre Termes}

% ---- Languages ----
\setdefaultlanguage[calendar=gregorian,hijricorrection=1]{arabic}
\setotherlanguage[variant=british]{english}

% ---- Hyperref (load after polyglossia/bidi) ----
\usepackage[hidelinks]{hyperref}
\hypersetup{
  pdftitle={السيرة الذاتية المهنية},
  pdfauthor={اسم المتقدم}
}

% ---- Section / list / table packages ----
\usepackage{titlesec}
\usepackage{enumitem}
\usepackage{tabularx}
\usepackage{booktabs}
\usepackage{xcolor}
\usepackage{fontawesome5}
\usepackage{setspace}

% ---- Colors ----
\definecolor{accent}{HTML}{1F3A5F}
\definecolor{rule}{HTML}{555555}

% ---- Section formatting (compact, colored, with rule) ----
% \titleformat* formats the starred (unnumbered) version of a section.
\titleformat*{\section}{\large\bfseries\color{accent}}
\titleformat*{\subsection}{\small\bfseries\color{accent}}
\titlespacing*{\section}{0pt}{0.7ex plus 0.2ex}{0.3ex plus 0.1ex}
\titlespacing*{\subsection}{0pt}{0.4ex plus 0.1ex}{0.2ex}

% Rule under each \section* heading (redefine \section to handle the star)
\makeatletter
\let\oldsection\section
\renewcommand{\section}{\@ifstar{\sectionstar}{\oldsection}}
\newcommand{\sectionstar}[1]{\oldsection*{#1}\vspace{-0.4em}{\color{rule}\hrule height 0.4pt}\vspace{0.25em}}
\makeatother

% ---- List spacing (very compact) ----
\setlist[itemize]{topsep=0pt, itemsep=0.05em, parsep=0pt, leftmargin=*}
\setlist[enumerate]{topsep=0pt, itemsep=0.05em, parsep=0pt, leftmargin=*}

% ---- Table column types ----
\newcolumntype{R}[1]{>{\raggedleft\arraybackslash}p{#1}}
\newcolumntype{C}[1]{>{\centering\arraybackslash}p{#1}}

% ---- Line spacing (very tight) ----
\setstretch{1.0}

\pagestyle{empty}

\begin{document}

% ============================================================
% HEADER
% ============================================================
\thispagestyle{empty}
\begin{center}
{\LARGE\bfseries\color{accent} اسم المتقدم\par}
\vspace{0.2em}
{\normalsize مهندس بيانات \,$\vert$\, محلل أنظمة\par}
\vspace{0.4em}
{\small
\faEnvelope~\href{mailto:name@example.com}{name@example.com} \quad
\faPhone~\LR{+90 555 000 0000} \quad
\faGlobe~\href{https://example.com}{example.com} \quad
\faMapMarker~\LR{Istanbul, Turkey}
\par}
\end{center}

\vspace{0.3em}

% ============================================================
% 1. الملخص
% ============================================================
\section*{الملخص}
مهندس بيانات بخبرة ست سنوات في بناء خطوط معالجة البيانات ونماذج التعلّم الآلي.
إتقان \LR{Python} و\LR{SQL} وأدوات البيانات الضخمة، مع سجل في تحسين الأنظمة
وزيادة الكفاءة التشغيلية.

% ============================================================
% 2. الخبرة العملية
% ============================================================
\section*{الخبرة العملية}

\textbf{مهندس بيانات أول} — شركة التقنية، إسطنبول \hfill \LR{2021–الآن}
\begin{itemize}
\item تصميم وتنفيذ خطوط معالجة البيانات باستخدام \LR{Apache Airflow} و\LR{Spark}.
\item بناء نماذج تنبؤ بالطلب أدت إلى خفض المخزون بنسبة \LR{18\%}.
\item قيادة فريق من أربعة مهندسين وتوثيق المعايير الداخلية.
\end{itemize}

\textbf{محلل بيانات} — شركة الحلول، أنقرة \hfill \LR{2018–2021}
\begin{itemize}
\item تحليل بيانات المبيعات وبناء لوحات تفاعلية بـ\LR{Tableau} و\LR{Power BI}.
\item أتمتة التقارير الأسبوعية باستخدام \LR{Python} ما وفّر عشر ساعات أسبوعياً.
\end{itemize}

% ============================================================
% 3. التعليم
% ============================================================
\section*{التعليم}
\vspace{-0.2em}
\begin{tabularx}{\textwidth}{@{}l X c@{}}
\toprule
\textbf{الدرجة} & \textbf{المؤسسة} & \textbf{السنة} \\
\midrule
ماجستير علوم البيانات & جامعة إسطنبول التقنية & \LR{2018} \\
بكالوريوس علوم حاسوب & جامعة الشرق الأوسط التقنية & \LR{2016} \\
\bottomrule
\end{tabularx}

% ============================================================
% 4. المهارات
% ============================================================
\section*{المهارات}
\textbf{تقنية:} \LR{Python}, \LR{SQL}, \LR{Spark}, \LR{Airflow}, \LR{Docker}, \LR{AWS}, \LR{Git}.\\
\textbf{لغوية:} العربية (اللغة الأم)، \LR{English} (إجادة تامة)، \LR{Turkish} (جيد).\\
\textbf{شخصية:} العمل ضمن فريق، حل المشكلات، إدارة الوقت، التواصل الفعّال.

% ============================================================
% 5. الشهادات
% ============================================================
\section*{الشهادات}
\begin{itemize}
\item مهندس بيانات معتمد من \LR{Google Cloud} — \LR{2022}.
\item شهادة \LR{AWS Solutions Architect} — \LR{2021}.
\item شهادة \LR{TensorFlow Developer} — \LR{2020}.
\end{itemize}

% ============================================================
% 6. اللغات
% ============================================================
\section*{اللغات}
العربية: اللغة الأم \quad\textbar\quad \LR{English}: إجادة تامة \quad\textbar\quad \LR{Turkish}: مستوى متوسط \quad\textbar\quad \LR{French}: مستوى أساسي.

\end{document}
